\begin{table}[ht]
\centering
\caption{Full leave-one-external-municipality-out sensitivity at the primary specification. For each omitted municipality, the directed strongly connected component, population and sex-specific kernels, residential distributions, and core-50 targets are re-estimated. The contrast compares the re-estimated regional female-minus-male gap with the unchanged Bogotá D.C. estimate. All entries are descriptive point estimates in FPT steps; the baseline contrasts are 0.520, 0.429, and 0.465 for employment, education, and health.}
\label{tab:s11_municipality_refit}
\begin{tabular}{llrr}
\toprule
Purpose & Omitted municipality & Refit contrast & Change from baseline \\
\midrule
Employment & Facatativá & 0.277 & -0.243 \\
Education & Facatativá & 0.080 & -0.348 \\
Health & Facatativá & 0.270 & -0.195 \\
Employment & Soacha & 0.389 & -0.131 \\
Education & Soacha & 0.560 & +0.131 \\
Health & Soacha & 0.347 & -0.118 \\
\bottomrule
\end{tabular}

\smallskip
\footnotesize Across all 20 omissions, the refit contrast ranges are [0.277, 0.565] for employment, [0.080, 0.564] for education, and [0.270, 0.552] for health. The corresponding minimum differences of normalized regional and district gaps are 0.052, 0.010, and 0.036. Facatativá has 229 sampled outbound households, 635 outbound trips, and 586 sampled residents in the regional network; Soacha has 1,449 households, 4,728 outbound trips, and 4,496 sampled residents. Each omission changes the regional model and the opportunity set; these differences are not causal effects.
\end{table}
